% QU-SLIDES-EXAMPLE.TEX
% Created 2022/05/01 by Joshua A. Marshall <joshua.marshall@queensu.ca>

\documentclass[fleqn]{qu-slides}

% ---------------------------------------------------------------------------
% INSTRUCTIONS
% 
% To change the background for a slide (and subsequent slides),
% invoke one of the following commands to switch backgrounds.
% 
% \SlideTitleBlue
% \SlideTitleRed
% \SlideTitleGold
% \SlideTitleGrey
% \SlideThemeWhite
% \SlideThemeBlue
% \SlideThemeGrey
% \SlideBlack
% \SlideBlue
% \SlideRed
% \SlideWhite
% \SlideGrey

% ---------------------------------------------------------------------------
% TEMPLATE OPTIONS

% Uncomment the line below if you want to remove the slide numbers
% \setbeamertemplate{footline}{~}

% To show the notes (for use, for example, with the SlidePilot app)
% \setbeameroption{show notes on second screen}

% ---------------------------------------------------------------------------
% PRESENTATION INFO

\title{The Presentation Title}
\subtitle{A Subtitle for the Presentation}
\author{Firstname M.~Lastname}
\institute{
    Department of Something \\
    Queen's University \\ 
    Kingston, Ontario, Canada
}
\date{\today}

% ---------------------------------------------------------------------------

\begin{document}

% ---------------------------------------------------------------------------
% TITLE SLIDES

% Default theme title slide (blue)
\SlideTitleBlue
\begin{frame}[plain]
    \titlepage
\end{frame}

% Alternate theme title slide (red)
\SlideTitleRed
\begin{frame}[plain]
    \titlepage
\end{frame}

% Alternate theme title slide (gold)
\SlideTitleGold
\begin{frame}[plain]
    \titlepage
\end{frame}

% Alternate theme title slide (grey)
\SlideTitleGrey
\begin{frame}[plain]
    \titlepage
\end{frame}

% ---------------------------------------------------------------------------
% OUTLINE SLIDE

\SlideThemeWhite
\begin{frame}
    \frametitle{Outline}
    \tableofcontents
\end{frame}

% ---------------------------------------------------------------------------
% NEW SECTION SLIDE

\section{Introduction}

\SlideThemeBlue
\begin{frame}[plain]
    \flushleft
    {\bf\LARGE \insertsection}
\end{frame}

% ---------------------------------------------------------------------------

\SlideThemeWhite
\begin{frame}
    \frametitle{Basic Theme Slide with Bullets}
    \begin{itemize}
        \item Itemized bullet number 1
        \begin{itemize}
            \item Sub-bullet number 1.1
            \item Sub-bullet number 1.2
        \end{itemize}
        \item Itemized bullet number 2
    \end{itemize}
\end{frame}

% This is how to add notes to the slide with the notes feature in Beamer
% \note[item]{This is note 1}
% \note[item]{This is note 2}

% ---------------------------------------------------------------------------

\SlideThemeWhite
\begin{frame}
    \frametitle{Colourful Slide with a Numbered List}
    \begin{enumerate}
        \item Enumerated item number 1
        \item Enumerated item number 2
    \end{enumerate}
\end{frame}

% ---------------------------------------------------------------------------
% NEW SECTION SLIDE

\section{Another Section}

\SlideThemeGrey
\begin{frame}[plain]
    \flushleft
    {\bf\LARGE \insertsection}
\end{frame}

% ---------------------------------------------------------------------------

\SlideWhite
\begin{frame}[fragile]
    \frametitle{Code with Syntax Highlighting}
    \begin{block}{Example Python Code Listing}
    \begin{minted}{python}
import numpy as np
x = np.zeros(5)
    \end{minted}
    \end{block}
\end{frame}

% ---------------------------------------------------------------------------

\SlideWhite
\begin{frame}[fragile]
    \frametitle{Example Theorem Environment}
    \begin{thm}
        The equilibrium point $\mbf{x}=\mbf{0}$ of $\dot{\mbf{x}}=\mbf{Fx}$ is asymptotically stable if and only if the eigenvalues $\lambda_{i}$ of $\mbf{F}$ satisfy $\textup{Re}\lambda_{i}<0$
    \end{thm}
\end{frame}

% ---------------------------------------------------------------------------

\SlideWhite
\begin{frame}[fragile]
    \frametitle{Example Alert Block}
    \begin{alertblock}{This is an Alert Block}
        Use this environment to draw attention to important stuff
    \end{alertblock}
\end{frame}

% ---------------------------------------------------------------------------

\SlideWhite
\begin{frame}
    \frametitle{Adding a Figure and Using a Blank Slide}
    \begin{columns}
        \column{0.5\textwidth}
        \begin{itemize}
            \item This is a standard bullet item
            \item[] This is another bullet item, but without the bullet
            \item Notice how to include a caption
        \end{itemize}
        \column{0.5\textwidth}
        \begin{figure}
            \flushleft
            \includegraphics[width=1.0\columnwidth]{logo/qu-logo-vertical-colour.pdf}
            \caption{This is the Queen's logo}
        \end{figure}
    \end{columns}
\end{frame}

% ---------------------------------------------------------------------------

\SlideWhite
\begin{frame}
    \frametitle{Here is a Table Example}
    \begin{table}
        \caption{This is an example of a table caption}
        \flushleft
        \begin{tabular}{lccc}
            \toprule
                \bf Small column & \multicolumn{3}{c}{\bf Big column} \\
            \midrule
                Grouped items & \multicolumn{3}{c}{Item 1} \\
            \cmidrule{2-4}
                & \multicolumn{3}{c}{Item 2} \\
            \midrule
                Usual row & Spam & Bacon & Eggs \\
            \bottomrule
        \end{tabular}
    \end{table}
\end{frame}

% ---------------------------------------------------------------------------

% TODO: Insert an example with a table

% ---------------------------------------------------------------------------

\SlideBlack
\begin{frame}
    \frametitle{This Slide is Black and Has Maths}
    \begin{columns}
        \column{0.5\textwidth}
        Try to integrate this function
        \begin{equation*}
            x_0=\int_0^\infty f_x(\tau)d\tau
        \end{equation*}
        where $\tau\in\Real$ is a variable
        \column{0.5\textwidth}
        \begin{thm}
            This is a very important theorem
        \end{thm}
    \end{columns}
\end{frame}

% ---------------------------------------------------------------------------

\SlideBlue 
\begin{frame}[plain]
    \centering
    \includegraphics[height=0.9\textheight]{logo/qu-logo-vertical-reverse.pdf}
\end{frame}

% ---------------------------------------------------------------------------

\SlideRed
\begin{frame}[plain]
    \centering
    \includegraphics[height=0.9\textheight]{logo/qu-logo-vertical-reverse.pdf}
\end{frame}

% ---------------------------------------------------------------------------

\SlideGrey
\begin{frame}[plain]
    \centering
    \includegraphics[height=0.9\textheight]{logo/qu-logo-vertical-colour.pdf}
\end{frame}

% ---------------------------------------------------------------------------

\SlideThemeGrey
\begin{frame}[plain]
    \begin{columns}
        \column{0.4\textwidth}
        \centering
        \includegraphics[width=\textwidth]{logo/qu-logo-vertical-colour.pdf}
        \column{0.6\textwidth}
        {\large\bf \insertauthor} \\
        \small \insertinstitute \\
        Tel: +1 613-533-6000 ext.\ 55555 \\
        Email: firstname.lastname@queensu.ca \\
        Web: www.queensu.ca
    \end{columns}
\end{frame}

% ---------------------------------------------------------------------------

\end{document}